\section{Conclusion}
\label{sec:conclusion}
A useful quantum reservoir must retain a readable record of the past while continuing to accept new information. We
have shown that the structure of dissipation can fundamentally change how this record survives. Independent decay
shortens memory as loss grows, whereas strong collective decay can prolong it by protecting information in dark states
through the quantum Zeno effect. The environment therefore provides a physical means of designing the memory of a
driven quantum system.

Longer storage also leaves more time for subsequent inputs to mix with the stored information. Our mixing law makes this accumulation quantitative. It sets a ceiling on the input strength:
the longer the memory, the more gently new information must be written, so that earlier inputs remain individually
recoverable. In the spin reservoirs studied, weakening the input raises the capacity to recall past inputs from about
15 to 35 units with exact readout, where no input is too weak.

Measurement adds a further requirement: the stored information must produce a signal strong enough to resolve.
Dark-state memories can remain protected for long times while becoming difficult to read. At strong loss a collective
channel with non-aligned phases keeps a readable memory where independent decay keeps none. Where the loss can be made
weak, independent decay recovers the most because its signals are stronger. The resources needed to retrieve information are therefore
part of the design of the memory itself.

Engineering the environment, the input, and the readout together thus provides a route from long-lived quantum
dynamics to memory that a readout can recover. Four rules summarize this design:

\begin{keybox}
\begin{enumerate}[label=(\arabic*),leftmargin=*,itemsep=2pt,topsep=0pt]
\item Collective decay can protect a memory. Stronger loss then makes it last longer, as long as local noise stays
  weak.
\item The input must be able to reach the protected memory. An input aligned with the decay channel hardly writes
  into it.
\item The longer a memory lasts, the more gently each input must be written. Otherwise new inputs mix with the old
  ones.
\item What can be read out depends on how strong a signal the inputs leave, so the measurement budget decides which
  design is best.
\end{enumerate}
\end{keybox}
